\documentclass[11pt]{article}
\usepackage[margin=1in]{geometry}
\usepackage[T1]{fontenc}
\usepackage[utf8]{inputenc}
\usepackage{amsmath,amssymb,amsthm}
\usepackage{enumitem}

\title{ICPC World Finals 2021\\G. Mosaic Browsing}
\author{}
\date{}

\begin{document}
\maketitle

\section*{Problem Summary}

We are given a large grid (the mosaic) and a smaller grid (the motif). The motif may contain zeroes,
which act as wildcards. We must output every top-left position where the motif matches the corresponding
subgrid of the mosaic.

The naive check of every position against every motif cell is too slow for $1000 \times 1000$ grids.

\section*{A One-Dimensional Identity}

For one aligned pair of cells, let
\[
    p = \text{motif value}, \qquad q = \text{mosaic value}.
\]

We want this cell to be valid exactly when either
\[
    p = 0
\]
or
\[
    p = q.
\]

The expression
\[
    p(q-p)^2
\]
does exactly that:
\begin{itemize}[leftmargin=*]
    \item if $p=0$, it is zero regardless of $q$;
    \item if $p \ne 0$, it is zero exactly when $q=p$;
    \item otherwise it is a positive integer.
\end{itemize}

So an alignment is valid if and only if
\[
    \sum p(q-p)^2 = 0
\]
over all aligned cells.

Expanding gives
\[
    \sum p q^2 - 2 \sum p^2 q + \sum p^3.
\]
The third sum depends only on the motif. The first two are cross-correlations, which can be computed by
FFT.

\section*{Turning the 2D Grid into 1D}

Flatten the mosaic in row-major order.

Flatten the motif in row-major order as well, but after each motif row insert
\[
    c_q - c_p
\]
zeroes, so that every motif row occupies exactly one full mosaic row in the flattened string.

Then placing the motif at top-left position $(r,c)$ in the 2D mosaic is exactly the same as aligning the
flattened padded motif with the flattened mosaic starting at index
\[
    (r-1)c_q + (c-1).
\]

This padding is the crucial trick: it prevents the end of one motif row from accidentally matching the
beginning of the next mosaic row.

\section*{Using Convolution}

Let the padded motif be $P$ and the flattened mosaic be $Q$. For each valid offset $s$, we need
\[
    \sum_i P_i (Q_{s+i} - P_i)^2 = 0.
\]

Expanding:
\[
    \sum_i P_i Q_{s+i}^2
    - 2 \sum_i P_i^2 Q_{s+i}
    + \sum_i P_i^3.
\]

We compute:
\begin{itemize}[leftmargin=*]
    \item the correlation of $P$ with $Q^2$;
    \item the correlation of $P^2$ with $Q$.
\end{itemize}

Each correlation is obtained by reversing the motif array and performing one ordinary convolution.

\section*{Algorithm}

\begin{enumerate}[leftmargin=*]
    \item Read the motif and mosaic.
    \item Flatten the mosaic into a 1D array $Q$, and also build $Q^2$.
    \item Flatten the motif with row padding into $P$, and also build $P^2$ and the constant
    \[
        C = \sum_i P_i^3.
    \]
    \item Reverse $P$ and $P^2$.
    \item Compute
    \[
        A = \text{conv}(\text{rev}(P), Q^2),
        \qquad
        B = \text{conv}(\text{rev}(P^2), Q).
    \]
    \item For every legal top-left position $(r,c)$, let
    \[
        s = (r-1)c_q + (c-1).
    \]
    The alignment score is
    \[
        A[s + |P| - 1] - 2B[s + |P| - 1] + C.
    \]
    This score is zero exactly for matches.
\end{enumerate}

\section*{Correctness Proof}

We prove that the algorithm outputs exactly all valid occurrences of the motif.

\paragraph{Lemma 1.}
For a single aligned cell with motif value $p$ and mosaic value $q$, the quantity
\[
    p(q-p)^2
\]
is zero if and only if the cell matches, meaning either $p=0$ or $p=q$.

\paragraph{Proof.}
If $p=0$, the value is clearly zero. If $p \ne 0$, then a square is zero only when $q-p=0$, i.e. when
$q=p$. In every other case it is strictly positive. \qed

\paragraph{Lemma 2.}
For one alignment of the motif against the mosaic, the total score
\[
    \sum_i P_i(Q_{s+i}-P_i)^2
\]
is zero if and only if the full alignment is a valid match.

\paragraph{Proof.}
By Lemma 1, every summand is a nonnegative integer and equals zero exactly when the corresponding cell
matches. Therefore the whole sum is zero exactly when every aligned cell matches. \qed

\paragraph{Lemma 3.}
The padded row-major flattening preserves exactly the legal 2D alignments.

\paragraph{Proof.}
Each motif row is padded to length $c_q$, so shifting by one row in the flattened motif corresponds to
shifting by exactly one full mosaic row. Therefore aligning the padded motif at flattened offset
\[
    (r-1)c_q + (c-1)
\]
compares motif cell $(i,j)$ with mosaic cell $(r+i-1,c+j-1)$, which is exactly the desired 2D alignment.
No cross-row wraparound can occur because of the inserted zero padding. \qed

\paragraph{Lemma 4.}
For every legal starting position, the algorithm computes the correct alignment score.

\paragraph{Proof.}
The term
\[
    \sum_i P_iQ_{s+i}^2
\]
is the correlation of $P$ with $Q^2$, and similarly
\[
    \sum_i P_i^2Q_{s+i}
\]
is the correlation of $P^2$ with $Q$. Reversing the motif arrays converts each correlation into an
ordinary convolution, and the index $s+|P|-1$ extracts the correct aligned value. Adding the constant
$\sum_i P_i^3$ gives exactly the expanded score from Lemma 2. \qed

\paragraph{Theorem.}
The algorithm outputs exactly all positions where the motif appears in the mosaic.

\paragraph{Proof.}
By Lemma 3, every legal 2D placement corresponds to one tested flattened offset. By Lemma 4, the
algorithm computes the exact score for that placement, and by Lemma 2 the score is zero exactly when the
placement is a valid match. Therefore the reported positions are exactly the motif occurrences. \qed

\section*{Complexity Analysis}

Let
\[
    N = r_q c_q, \qquad M = (r_p-1)c_q + c_p.
\]
We perform two FFT-based convolutions of size $O(N+M)$, so the total running time is
\[
    O((N+M)\log(N+M)).
\]
The memory usage is linear in the FFT size.

\section*{Implementation Notes}

\begin{itemize}[leftmargin=*]
    \item The score is always a nonnegative integer, so after FFT rounding errors it is safe to accept a
    match when the computed value is very close to zero.
    \item If $r_p > r_q$ or $c_p > c_q$, there are no legal placements at all.
\end{itemize}

\end{document}
